%% ══════════════════════════════════════════════════════════════
%%  第1章 绪论
%% ══════════════════════════════════════════════════════════════

\chapter{绪论}

\label{chap:intro}

%% ── 1.1 研究背景与意义 ─────────────────────────────────────

\section{研究背景与意义}

随着人口老龄化加剧和劳动力成本上升，协作机器人（Cobot）在制造业、服务业、医疗康复等领域的应用需求快速增长。与传统工业机器人相比，协作机器人强调人机共融、安全易用和快速部署，这对机器人系统的体积、功耗、成本和实时性提出了更严苛的要求。

近年来，大模型与机器人学习技术的融合成为研究热点。视觉-语言-动作模型（Vision-Language-Action Model，VLA）证明了大数据预训练范式在机器人控制中的有效性。与此同时，模仿学习策略如 Action Chunking with Transformers（ACT）也在低成本硬件上展现出良好的任务泛化能力。然而，将这些 AI 模型部署到资源受限的嵌入式平台（如树莓派等 ARM 设备）面临严重的算力瓶颈。树莓派 5 搭载 Cortex-A76 处理器，无 GPU 加速，NEON SIMD 能力有限，使得感知-推理-执行流水线的端到端延迟成为制约实时控制的关键因素。

因此，对机器人学习框架进行系统级性能剖析，识别感知、推理、执行各层的性能瓶颈，并针对性优化，对于推动协作机器人在边缘设备上的实际部署具有重要的理论意义和工程价值。

%% ── 1.2 国内外研究现状 ─────────────────────────────────────

\section{国内外研究现状}

机器人学习领域已形成以模仿学习和强化学习为主的两条技术路线。模仿学习方面，ACT 策略网络通过动作分块和时间序列 Transformer 实现了低延迟、高精度的双臂 manipulation 控制；Diffusion Policy 将动作生成建模为扩散过程，在复杂多模态任务上表现优异。VLA 模型方面，Google 的 RT-1/RT-2 系列首次验证了大规模机器人数据训练 Transformer 的可行性，OpenVLA 和 π0 则将 VLA 范式推向开源化和通用化。

在开源机器人框架层面，HuggingFace 的 LeRobot 提供了统一的数据集格式、策略接口和训练流程，成为低コスト机器人学习的主流选择。Mobile ALOHA 结合仿生移动机械臂与模仿学习，以低成本硬件实现了复杂的全身 manipulation 任务。

在软件架构层面，ROS2 作为机器人操作系统的事实标准，提供了分布式通信、生命周期管理和硬件抽象等基础设施。然而其默认 DDS 通信延迟较高，不适合实时控制场景。轻量级通信库 LCM 曾在实时性要求较高的场景中被广泛采用。此外，移动端推理框架（ONNX Runtime、TFLite、PyTorch Mobile）和 ARM NEON 指令优化也为边缘部署提供了多条技术路径。

在开源硬件与低成本本体方面，多个平台为机器人学习研究提供了可复现的硬件基础。ROBOTIS 的 OpenManipulator 系列提供配套 ROS2 的开源机械臂，硬件图纸与代码全开源，适合教学和原型开发。TurtleBot 系列作为开源移动机器人的事实标准，最新版本 TurtleBot 4 与 ROS2 生态深度集成。UC Berkeley HybridRobotics 实验室开发的 Berkeley Humanoid Lite 采用 3D 打印方案，为小型人形机器人的开源化提供了新思路。

本课题基于 LeRobot 框架，以树莓派 5 为部署平台，对 ACT 策略网络的感知-推理-执行流水线进行系统级性能剖析，涵盖内核层、框架层、AI 算子层和应用层四个维度，旨在为协作机器人在边缘嵌入式设备上的高效部署提供理论依据和优化方案。

%% ── 1.3 本文主要研究内容 ─────────────────────────────────

\section{本文主要研究内容}

本文围绕"基于树莓派 5 的 LeRobot 机器人学习框架性能剖析"这一课题，开展以下研究工作：

\begin{enumerate}
  \item \textbf{搭建基于 LeRobot 的机器人感知-推理-执行实验平台：}在树莓派 5 上部署 LeRobot 框架，以 SO-101 仿生机械臂为执行本体，构建完整的感知-推理-执行控制流水线。

  \item \textbf{建立四层 profiling 分析模型：}从内核层（ftrace 函数追踪）、框架层（ROS2 通信延迟）、AI 算子层（torch.profiler）和应用层（端到端延迟）四个维度，建立系统级性能剖析的方法论。

  \item \textbf{识别流水线瓶颈并量化各层开销：}通过分层测量，精确定位感知、推理、执行各层的性能瓶颈，量化各层的延迟贡献。

  \item \textbf{探索 ARM 嵌入式平台上的推理优化路径：}分析 INT8/FP16 量化、NEON SIMD 加速和多线程并行等优化手段在树莓派 5 上的可行性与效果。
\end{enumerate}

%% ── 1.4 论文结构安排 ─────────────────────────────────────

\section{论文结构安排}

本文共分为六章，各章内容安排如下：

\textbf{第 1 章} 绪论。阐述研究背景与意义，分析国内外研究现状，明确本文的主要研究内容与结构安排。

\textbf{第 2 章} 相关技术与理论基础。介绍 LeRobot 框架、ACT 策略网络、VLA 模型等核心技术的原理，以及性能剖析的方法论工具。

\textbf{第 3 章} 实验平台搭建。详细描述树莓派 5 硬件环境、LeRobot 框架部署、SO-101 机械臂集成和控制流水线的构建过程。

\textbf{第 4 章} 性能剖析与瓶颈分析。基于四层 profiling 模型，对感知-推理-执行流水线进行系统级性能测试，识别各层瓶颈。

\textbf{第 5 章} 优化方案与实验验证。针对识别的瓶颈，提出并验证适用于 ARM 嵌入式平台的推理优化方案。

\textbf{第 6 章} 总结与展望。总结本文的研究工作，展望未来的研究方向。

\endinput
